\Chapitre{Étiquetage}

\section{Modules}

Sur les modules, les inscriptions sont faites sur chacune des faces de raccord

Les inscriptions contiennent :

\begin{itemize}
    \item le propriétaire
    \item le n° du module
    \item le n° de sa caisse de rangement
    \item son emplacement de rangement dans la caisse, généralement l'étage
\end{itemize}

% TODO prévoir des couleurs pour chaque indications

\section{Caisses de modules}

Sur les caisses de module, les inscriptions sont faites au minimum sur une des faces et sur le couvercle.

Les inscriptions contiennent :

\begin{itemize}
    \item le propriétaire
    \item le n° de la caisse
    \item la liste des modules disposé de manière similaire à l'emplacement des modules dans la caisse
    \item le poids de la caisse pleine en kg
    \item la hauteur de la caisse
\end{itemize}

Les inscriptions peuvent également contenir :

\begin{itemize}
    \item le nom de la gare ou du groupe de module pour lequel la caisse est nécessaire
    \item la liste des caisses qui sont nécessaire de prendre avec
    \item le lieu de stockage de la caisse
    \item les information de contact du propriétaire
\end{itemize}

\section{caisse de matériel}

Sur les caisses de matériel empilables, les inscriptions sont faites au minimum sur le devant et sur un des côté.

Les inscriptions contiennent :

\begin{itemize}
  \item le propriétaire
  \item le n° de la caisse
  \item la liste ou la désignation générale du contenu
  \item le poids de la caisse pleine en kg
\end{itemize}

Les inscriptions peuvent également contenir :

\begin{itemize}
  \item le nom de la gare ou du groupe de module pour lequel la caisse est nécessaire
  \item la liste des caisses qui sont nécessaire de prendre avec
  \item le lieu de stockage de la caisse
  \item les information de contact du propriétaire
\end{itemize}
